\section{Big Data}
\label{sec:ch06-big-data}
\begin{definitionbox}{}
যে ডেটার পরিমাণ, উৎপাদনের গতি বা বৈচিত্র্যের কারণে প্রচলিত ব্যবস্থায় কাঙ্ক্ষিত সময়ে কার্যকরভাবে সংরক্ষণ, প্রক্রিয়াকরণ ও বিশ্লেষণ কঠিন হয়, তাকে বিগ ডেটা বলে।
\end{definitionbox}


একটি দোকানের দৈনিক বিক্রয়ের সারণি ছোট হতে পারে। কিন্তু বহু শাখার বিক্রয়, ওয়েবসাইটে ক্রেতার ক্লিক, পণ্যের ছবি ও ক্রেতার মতামত একত্র করলে ডেটার পরিমাণ বেড়ে যায়। এর সঙ্গে প্রতি মুহূর্তে পণ্য পৌঁছে দেওয়ার সর্বশেষ অবস্থাও যোগ হতে থাকে। তখন শুধু পরিমাণ নয়, তথ্য আসার গতি ও বিভিন্ন ধরনের তথ্যও সামলাতে হয়।

শুধু ফাইল বড় হলেই বা নির্দিষ্ট সংখ্যক গিগাবাইট (Gigabyte বা GB) পার হলেই সব ক্ষেত্রে Big Data হয়—এমন স্থির সীমা নেই। ব্যবস্থার সক্ষমতা ও কাজের সময়সীমাও বিবেচ্য। পরীক্ষার ফলের সুসংগঠিত টেবিল, JSON বিন্যাসের লগ এবং ছবি বা অবাধে লেখা মতামত ভিন্ন ধরনের তথ্য। বৃহৎ বিশ্লেষণে এগুলো একসঙ্গে কাজে লাগতে পারে।


\begin{flushleft}\small
\begin{tabularx}{\linewidth}{|p{27mm}|Y|Y|}\hline\chSixHeader
গঠন & চেনার উপায় & পরিচিত উদাহরণ\\\hline
Structured & নির্ধারিত ক্ষেত্র ও সারির বিন্যাস & Roll, Name, GPA-এর সারণি\\\hline
Semi-structured & key/tag দিয়ে কিছু সংগঠন; ক্ষেত্র বদলাতে পারে & JSON-এ পণ্যের বিবরণ\\\hline
Unstructured & বিশ্লেষণের তথ্য পূর্বনির্ধারিত সারি--কলামে সাজানো নয় & মুক্ত review, ছবি, audio\\\hline
\end{tabularx}
\end{flushleft}
Unstructured ডেটা মানে ফাইলের নিজস্ব ফরম্যাট নেই—এমন নয়। ছবির ফরম্যাট থাকলেও তার বিষয়বস্তু শিক্ষার্থীর সারণির মতো সাজানো নয়। আবার RDBMS-এ ছবির মতো বাইনারি ডেটা বা JSON রাখার সুবিধা থাকতে পারে। তাই ভিডিও আছে বলেই RDBMS বাদ দিতে হবে—এ সিদ্ধান্ত ঠিক নয়; সংরক্ষণ ও বিশ্লেষণের প্রয়োজন বিচার করতে হয়।

\subsection{Big Data-এর 5V বৈশিষ্ট্য}
Volume, Velocity ও Variety মৌলিক তিনটি দিক। ডেটার বিশ্বাসযোগ্যতা ও উপযোগিতা বোঝাতে Veracity ও Value-সহ পাঁচটি V-ও আলোচনা করা হয়। এগুলো বোঝার কাঠামো; সব বই বা ব্যবস্থায় বৈশিষ্ট্যের সংখ্যা ঠিক পাঁচটিই হতে হবে না।
\begin{enumerate}
\item \textbf{Volume:} ডেটার বিপুল পরিমাণ বোঝায়। উদাহরণ: বহু বছরের অর্ডার, click ও ছবির বড় সংগ্রহ।
\item \textbf{Velocity:} তথ্য আসা ও সামলানোর গতি বোঝায়। উদাহরণ: প্রতি মুহূর্তে নতুন অর্ডার ও delivery update।
\item \textbf{Variety:} গঠন ও ধরনের বৈচিত্র্য বোঝায়। উদাহরণ: বিক্রয়ের সারণি, JSON log, ছবি ও review।
\item \textbf{Veracity:} তথ্যের যথার্থতা ও বিশ্বাসযোগ্যতা বোঝায়। উদাহরণ: ভুয়া review, duplicate order ও ভুল ঠিকানা শনাক্ত করা।
\item \textbf{Value:} বিশ্লেষণের কার্যকর উপযোগিতা বোঝায়। উদাহরণ: চাহিদা বুঝে মজুত পরিকল্পনা ও delivery উন্নতি।
\end{enumerate}
অনেক মতামত থাকা Volume, দ্রুত নতুন মতামত আসা Velocity এবং লেখা-ছবি-ভিডিও একসঙ্গে থাকা Variety বোঝাতে পারে। মতামত সত্য কি না, তা Veracity-এর বিষয়; সেই তথ্য ব্যবহার করে সেবার উন্নতি করা Value বোঝায়। একই ঘটনার একাধিক দিক থাকতে পারে, তাই প্রশ্নে কোন বৈশিষ্ট্যের ওপর জোর দেওয়া হয়েছে তা পড়তে হবে।


\subsubsection{ওয়েবসাইটের ব্যবহার বিশ্লেষণে Big Data}
\begin{enumerate}
\item \textbf{ঘটনার তথ্য সংগ্রহ:} একটি অনলাইন দোকানে প্রতিদিন অর্ডার, পণ্য দেখার তথ্য, অনুসন্ধানে লেখা শব্দ এবং ছবিসহ ক্রেতার মতামত জমা হয়। প্রয়োজনীয় ঘটনার তথ্য সময় ও পণ্যের পরিচয়সহ সংগ্রহ করতে হয়।
\item \textbf{চাহিদা ও সেবার বিশ্লেষণ:} শুধু ফাইল সংরক্ষণ করলে কোন পণ্যের চাহিদা বাড়ছে বা কোথায় সেবা ধীর হচ্ছে, তা বোঝা যায় না। ভুল ও পুনরাবৃত্তি যাচাই করে উপযুক্ত হিসাব ও বিশ্লেষণ করতে হয়।
\item \textbf{কাজ ভাগ করে সামলানো:} ব্যবস্থার সামর্থ্য কম হলে কাজ ভাগ করে সামলানোর প্রয়োজন হতে পারে।
\item \textbf{অন্য ক্ষেত্র:} সামাজিক যোগাযোগের ব্যবহার-প্রবণতা, পরিবহনের চাপ এবং হাসপাতালের সেবার চাহিদা বিশ্লেষণেও একই ধরনের প্রয়োজন দেখা যায়।
\end{enumerate}
\textbf{গণনার শর্ত:} একই ব্যক্তি বারবার কোনো পণ্য দেখলে সেই একজনকে একাধিক ক্রেতা হিসেবে গণনা করা যাবে না। সংখ্যার পাশাপাশি সেই সংখ্যা কোন ঘটনা বোঝাচ্ছে, তাও বুঝতে হবে।

\subsubsection{5V বৈশিষ্ট্যের সমন্বিত তাৎপর্য}
\begin{enumerate}
\item \textbf{পরিমাণ, গতি ও বৈচিত্র্য:} Volume, Velocity ও Variety ডেটার পরিমাণ, আগমনের গতি ও বৈচিত্র্য সামলানোর প্রয়োজন বোঝায়।
\item \textbf{বিশ্বাসযোগ্যতা ও উপযোগিতা:} Veracity তথ্যের নির্ভরযোগ্যতা এবং Value সেই তথ্যের উপকারিতা বোঝায়।
\item \textbf{সমন্বিত ব্যবস্থাপনা:} অনেক সেন্সরের দ্রুত আসা পাঠ সংরক্ষণ করলেই কাজ শেষ হয় না। ত্রুটিপূর্ণ পাঠ শনাক্ত করে সময়মতো সিদ্ধান্ত নেওয়াও প্রয়োজন। বিগ ডেটা বোঝার ক্ষেত্রে তথ্যের আকারের সঙ্গে ব্যবস্থাপনা ও বিশ্লেষণের প্রয়োজন বিবেচনা করতে হয়।
\end{enumerate}
নিচের প্রবাহে সংগ্রহ থেকে সিদ্ধান্ত পর্যন্ত কাজগুলো দেখানো হয়েছে।


\subsection{উৎস থেকে সিদ্ধান্ত}
\begin{flushleft}\centering\captionsetup{justification=centering}
\begin{tikzpicture}[node distance=4mm,every node/.style={booknode,font=\small,text width=25mm,minimum height=11mm}]
\node (a) {Sources};
\node[right=of a] (b) {Collect / Store};
\node[right=of b] (c) {Check / Analyse};
\node[right=of c] (d) {Decision};
\draw[bookarrow](a)--(b);\draw[bookarrow](b)--(c);\draw[bookarrow](c)--(d);
\end{tikzpicture}
\end{flushleft}
অনলাইন লেনদেন, যন্ত্রের সেন্সর, বিভিন্ন সেবার লগ ও সামাজিক যোগাযোগের প্রকাশযোগ্য তথ্য সম্ভাব্য উৎস। সংগ্রহের পরে ভুল, অসম্পূর্ণতা ও পুনরাবৃত্তি যাচাই করে প্রশ্ন অনুযায়ী বিশ্লেষণ করতে হয়। যেমন কোন এলাকায় কোন পণ্যের চাহিদা বাড়ছে জানলে সেখানে মজুত বাড়ানোর সিদ্ধান্ত নেওয়া যায়। শুধু বিপুল তথ্য জমিয়ে রাখলেই সিদ্ধান্ত সঠিক হয় না।

\subsection{ব্যবহার}
\begin{enumerate}
\item \textbf{পরিবহন:} যাত্রীর চাপ বিশ্লেষণ করে সেবা পরিকল্পনায় সহায়তা করে।
\item \textbf{স্বাস্থ্যসেবা:} হাসপাতালের সেবার চাহিদা ও রোগের প্রবণতা বোঝায়।
\item \textbf{কৃষি:} আবহাওয়া ও উৎপাদনের তথ্য বিশ্লেষণে কাজে লাগে।
\end{enumerate}
\subsubsection{সুবিধা}
\begin{enumerate}
\item \textbf{প্রবণতা শনাক্ত:} অনেক উৎসের তথ্য থেকে পরিবর্তনের ধারা বোঝা যায়।
\item \textbf{প্রমাণভিত্তিক পরিকল্পনা:} বিশ্লেষণের ফল দিয়ে সিদ্ধান্ত গ্রহণে সহায়তা করে।
\end{enumerate}
\subsubsection{সীমাবদ্ধতা}
\begin{enumerate}
\item \textbf{ভুল সিদ্ধান্তের ঝুঁকি:} অসম্পূর্ণ বা পক্ষপাতপূর্ণ নমুনা বিভ্রান্তিকর ফল দিতে পারে।
\item \textbf{ব্যয় ও দক্ষতা:} তথ্য মেলানো, সংরক্ষণ ও বিশ্লেষণে খরচ এবং দক্ষ লোকবল লাগে।
\item \textbf{গোপনীয়তা:} ব্যক্তিগত তথ্য ব্যবহারে প্রয়োজনীয় অনুমতি, সীমিত অধিকার ও পরিচয় সুরক্ষা প্রয়োজন।
\end{enumerate}
বিশ্লেষণের ফল বাস্তব পরিস্থিতির সঙ্গে মিলিয়ে দেখতে হয়; জনসেবার তথ্য বিশ্লেষণ করতে সবার ব্যক্তিগত বিবরণ প্রকাশ করা প্রয়োজন নেই।

\subsection{Big Data, সাধারণ Data ও NoSQL}
\begin{flushleft}\small
\begin{tabularx}{\linewidth}{|p{29mm}|Y|}\hline\chSixHeader
ধারণা & মূল বিবেচনা\\\hline
ছোট ফলের ডেটাবেজ & নির্দিষ্ট কয়েকটি সারণি ও পরিচিত কুয়েরি দিয়ে কাজ সম্ভব\\\hline
Big Data & পরিমাণ, গতি ও বৈচিত্র্যের সঙ্গে বিশ্লেষণের চাহিদা সামলানো\\\hline
NoSQL & ডেটা সংগঠন ও পরিচালনার বিকল্প মডেলভিত্তিক ব্যবস্থা\\\hline
\end{tabularx}
\end{flushleft}
Big Data একটি ডেটা ও প্রক্রিয়াকরণের প্রসঙ্গ; NoSQL ডেটাবেজ ব্যবস্থার শ্রেণি। Big Data-র কাজে NoSQL সহায়ক হতে পারে, কিন্তু সব Big Data অবশ্যই NoSQL-এ রাখতে হবে না। প্রচলিত রিলেশনাল ব্যবস্থা, অন্য উপযোগী প্রযুক্তি বা একাধিক ব্যবস্থা প্রয়োজন অনুযায়ী একসঙ্গে ব্যবহার করা যায়।
